% -*- coding: utf-8 -*-
% Marcas para la revisión manual de los apuntes. Instrucciones en COMO_REVISAR.md.
%
%   \comenta{texto}            comentario u observación
%   \duda{texto}               pregunta o duda para resolver
%   \pendiente{texto}          tarea pendiente
%   \corrige{original}{nuevo}  propuesta de sustitución
%   \borra{texto}              propuesta de eliminación
%   \inserta{texto}            propuesta de inserción
%   \respuesta{texto}          respuesta de Claude a una duda o comentario
%
% \revisiontrue  -> las marcas se ven en el PDF (con color y numeradas).
% \revisionfalse -> PDF limpio: se ocultan comentarios, dudas y pendientes;
%                   \corrige muestra la versión nueva, \borra elimina el texto
%                   e \inserta lo muestra sin marcar.

\usepackage[normalem]{ulem}   % \sout (tachado) y \uline (subrayado); normalem mantiene \emph en cursiva

\newif\ifrevision
\revisiontrue

\definecolor{revComenta}{RGB}{0,90,170}
\definecolor{revDuda}{RGB}{200,90,0}
\definecolor{revPendiente}{RGB}{140,0,140}
\definecolor{revBorra}{RGB}{190,0,0}
\definecolor{revInserta}{RGB}{0,130,60}
\definecolor{revRespuesta}{RGB}{90,90,90}

\newcounter{marcarev}

% Texto en modo texto aunque la marca se use dentro de una fórmula.
\newcommand{\revtexto}[1]{\ifmmode\text{#1}\else#1\fi}

% Marca numerada: #1 color, #2 tipo, #3 texto.
\newcommand{\revmarca}[3]{%
  \stepcounter{marcarev}%
  \revtexto{\textcolor{#1}{\sffamily\small[\textbf{#2~\arabic{marcarev}:} #3]}}}

\DeclareRobustCommand{\comenta}[1]{\ifrevision\revmarca{revComenta}{Nota}{#1}\fi}
\DeclareRobustCommand{\duda}[1]{\ifrevision\revmarca{revDuda}{Duda}{#1}\fi}
\DeclareRobustCommand{\pendiente}[1]{\ifrevision\revmarca{revPendiente}{Pendiente}{#1}\fi}
\DeclareRobustCommand{\respuesta}[1]{\ifrevision\revmarca{revRespuesta}{Respuesta}{#1}\fi}

% Tachado y subrayado. Dentro de una fórmula (por ejemplo, en una línea de
% una derivación Fitch) el argumento se trata como fórmula.
\newcommand{\revtachado}[1]{\ifmmode\text{\textcolor{revBorra}{\sout{$#1$}}}\else\textcolor{revBorra}{\sout{#1}}\fi}
\newcommand{\revnuevo}[1]{\ifmmode\text{\textcolor{revInserta}{\uline{$#1$}}}\else\textcolor{revInserta}{\uline{#1}}\fi}

\DeclareRobustCommand{\borra}[1]{\ifrevision\revtachado{#1}\fi}
\DeclareRobustCommand{\inserta}[1]{\ifrevision\revnuevo{#1}\else#1\fi}
\DeclareRobustCommand{\corrige}[2]{\ifrevision\revtachado{#1}\,\revnuevo{#2}\else#2\fi}

% En los marcadores del PDF (hyperref) las marcas no se pueden componer:
% se ignoran los comentarios y se usa la versión propuesta.
\pdfstringdefDisableCommands{%
  \def\comenta#1{}\def\duda#1{}\def\pendiente#1{}\def\respuesta#1{}%
  \def\borra#1{}\def\inserta#1{#1}\def\corrige#1#2{#2}}

% El índice se compone sin marcas (versión limpia), para que no se
% dupliquen ni alteren la numeración de las marcas del texto.
\let\revtableofcontentsoriginal\tableofcontents
\renewcommand{\tableofcontents}{\begingroup\revisionfalse\revtableofcontentsoriginal\endgroup}
